%
%  chapter3.tex  —  Resultados Experimentais e Conclusões
%
\chapter{Resultados Experimentais e Conclusões}
\label{cha:resultados}

\section{Configuração Experimental}
\label{sec:config_exp}

Os testes foram realizados num MacBook Pro com processador Apple M-series, sistema operativo macOS, SWI-Prolog 10.0.2 e Python 3.14.3. Foram utilizados quatro puzzles clássicos de dificuldade crescente e um puzzle Killer Sudoku (fácil):

\begin{itemize}
    \item \textbf{Fácil} (\texttt{classic\_easy\_01}) — 51 células vazias;
    \item \textbf{Médio} (\texttt{classic\_medium\_01}) — 51 células vazias;
    \item \textbf{Difícil} (\texttt{classic\_hard\_01}) — puzzle AI Escargot (17 \emph{givens});
    \item \textbf{Especialista} (\texttt{classic\_expert\_01}) — 51 células vazias;
    \item \textbf{Killer Fácil} (\texttt{killer\_easy\_01}) — 81 células vazias, 27 gaiolas.
\end{itemize}

Os solucionadores SA e GA foram executados 3 vezes por puzzle (natureza estocástica). O limite de tempo de Prolog foi de 120 segundos por chamada. O SA usou no máximo $300\,000$ iterações. O GA usou população de 200 indivíduos e no máximo $10\,000$ gerações.

\section{Resultados dos Solucionadores Prolog}
\label{sec:resultados_prolog}

A Tabela~\ref{tab:prolog_classic} apresenta os resultados do DFID e A* para os puzzles clássicos.

\begin{table}[ht]
    \caption{Resultados do DFID e A* nos puzzles clássicos.}
    \label{tab:prolog_classic}
    \centering
    \begin{tabular}{lrrrr}
        \toprule
        \textbf{Puzzle} & \multicolumn{2}{c}{\textbf{DFID}} & \multicolumn{2}{c}{\textbf{A*}} \\
        \cmidrule(r){2-3} \cmidrule(r){4-5}
        & Nós & Tempo (s) & Nós & Tempo (s) \\
        \midrule
        Fácil    &    51 & 0.163 &    51 & 0.173 \\
        Médio    &    49 & 0.153 &    49 & 0.160 \\
        Difícil  &   481 & 2.671 &  2276 & 13.314 \\
        Especialista &  846 & 3.667 &   216 & 1.024 \\
        \bottomrule
    \end{tabular}
\end{table}

\paragraph{Análise DFID.}
O DFID com MRV expande poucos nós em todos os puzzles. O MRV reduz drasticamente o factor de ramificação ao escolher sempre a célula com menos candidatos válidos. Em puzzles bem constrangidos (fácil, médio), o número de nós expandidos é inferior a 55, o que confirma que a pesquisa é efectivamente uma DFS com verificação de restrições que raramente precisa de retroceder. No puzzle Difícil (AI Escargot, com apenas 17 \emph{givens}), o DFID expande 481 nós — o maior esforço de pesquisa entre os casos testados com este algoritmo.

\paragraph{Análise A*.}
O A* expande exactamente o mesmo número de nós que o DFID nos puzzles Fácil e Médio. No puzzle Difícil, o A* expande 2276 nós (vs. 481 do DFID), porque a heurística $h = $ células vazias é constante ao mesmo nível de profundidade, fazendo o A* comportar-se como BFS. A fila de prioridade cresce rapidamente, aumentando o overhead por nó. No puzzle Especialista, o A* expande apenas 216 nós vs. 846 do DFID, possivelmente devido a diferenças na ordenação da fila de prioridade em casos de empate de prioridade.

\paragraph{Resultados Killer Sudoku — Prolog.}
A Tabela~\ref{tab:prolog_killer} mostra os resultados para o puzzle Killer Fácil.

\begin{table}[ht]
    \caption{Resultados do DFID e A* no Killer Sudoku (fácil).}
    \label{tab:prolog_killer}
    \centering
    \begin{tabular}{lrr}
        \toprule
        \textbf{Algoritmo} & \textbf{Nós Expandidos} & \textbf{Tempo (s)} \\
        \midrule
        DFID Killer & 51 & 0.188 \\
        A* Killer   & 51 & 0.176 \\
        \bottomrule
    \end{tabular}
\end{table}

O puzzle Killer Fácil (sem \emph{givens} iniciais) é resolvido com apenas 51 nós por ambos os algoritmos. A poda por \texttt{partial\_cage\_ok} é muito eficaz: as restrições de gaiola eliminam candidatos logo nas primeiras atribuições, reduzindo o espaço de busca de forma semelhante aos \emph{givens} de um puzzle clássico.

\section{Resultados do Arrefecimento Simulado (SA)}
\label{sec:resultados_sa}

A Tabela~\ref{tab:sa} apresenta os resultados do SA nas 3 execuções por puzzle.

\begin{table}[ht]
    \caption{Resultados do SA (3 execuções por puzzle; limite: 300\,000 iterações).}
    \label{tab:sa}
    \centering
    \begin{tabular}{llrrr}
        \toprule
        \textbf{Puzzle} & \textbf{Resultado} & \textbf{Execuções Resolvidas} & \textbf{Iter. (média)} & \textbf{Tempo (s)} \\
        \midrule
        Fácil       & Resolvido   & 3/3 & $\approx$17\,644 & 0.06--0.10 \\
        Médio       & Resolvido   & 3/3 & $\approx$44\,223 & 0.08--0.38 \\
        Difícil     & Não resolvido & 0/3 & 300\,000 & $\approx$1.21 \\
        Especialista & Parcial    & 1/3 & $\approx$95\,647 & 0.38 \\
        Killer Fácil & Resolvido  & 3/3 & $\approx$119\,560 & 0.87--1.58 \\
        \bottomrule
    \end{tabular}
\end{table}

O SA resolve de forma fiável os puzzles Fácil e Médio em poucas dezenas de milissegundos. O puzzle Difícil (AI Escargot) constitui o caso mais desafiante: em todas as execuções, o algoritmo fica bloqueado num estado com custo 2, que corresponde a duas violações residuais em linhas/colunas distintas, cujas células envolvidas pertencem a caixas diferentes. O operador de troca intra-caixa não consegue resolver este conflito sem aceitar temporariamente estados de custo superior, o que a temperatura baixa torna improvável. O puzzle Especialista é resolvido em 1 de cada 3 execuções, demonstrando a variabilidade estocástica do algoritmo.

O puzzle Killer Fácil é resolvido em todas as execuções, apesar de não ter \emph{givens} iniciais, porque as restrições de gaiola reduzem substancialmente o espaço efectivo de soluções.

\section{Resultados do Algoritmo Genético (GA)}
\label{sec:resultados_ga}

A Tabela~\ref{tab:ga} apresenta os resultados do GA.

\begin{table}[ht]
    \caption{Resultados do GA (3 execuções; pop.=200, máx. 10\,000 gerações, 15--26 s/execução).}
    \label{tab:ga}
    \centering
    \begin{tabular}{llrr}
        \toprule
        \textbf{Puzzle} & \textbf{Resultado} & \textbf{Execuções Resolvidas} & \textbf{Custo Mínimo} \\
        \midrule
        Fácil        & Não resolvido & 0/3 & 6 \\
        Médio        & Não resolvido & 0/3 & 2 \\
        Difícil      & Não resolvido & 0/3 & 4 \\
        Especialista & Não resolvido & 0/3 & 4 \\
        Killer Fácil & Não resolvido & 0/3 & 15 \\
        \bottomrule
    \end{tabular}
\end{table}

O GA não encontrou uma solução completa em nenhuma das execuções realizadas. O principal problema é a \emph{convergência prematura}: a diversidade genética da população colapsa rapidamente, e o GA queda num mínimo local sem conseguir sair, mesmo com a reinicialização por estagnação. O cruzamento preservador de caixas, apesar de manter as restrições de caixa, não introduz diversidade suficiente para explorar espaços de soluções de dimensão elevada.

O custo mínimo atingido pelo GA é consistentemente superior ao do SA nos mesmos puzzles. No puzzle Killer Fácil, o custo mínimo de 15 reflecte que o GA raramente satisfaz as restrições de gaiola, que exigem uma combinação muito específica de valores.

\section{Análise Comparativa}
\label{sec:analise_comp}

\begin{table}[ht]
    \caption{Comparação dos quatro solucionadores nos puzzles clássicos.}
    \label{tab:comparacao}
    \centering
    \begin{tabular}{lcccc}
        \toprule
        \textbf{Puzzle} & \textbf{DFID} & \textbf{A*} & \textbf{SA} & \textbf{GA} \\
        \midrule
        Fácil       & \checkmark\ 0.16s & \checkmark\ 0.17s & \checkmark\ 0.06--0.10s & $\times$\ custo 6 \\
        Médio       & \checkmark\ 0.15s & \checkmark\ 0.16s & \checkmark\ 0.08--0.38s & $\times$\ custo 2 \\
        Difícil     & \checkmark\ 2.67s & \checkmark\ 13.3s & $\times$\ custo 2     & $\times$\ custo 4 \\
        Especialista & \checkmark\ 3.67s & \checkmark\ 1.02s & 1/3\ resolvido        & $\times$\ custo 4 \\
        \bottomrule
    \end{tabular}
\end{table}

Os algoritmos de pesquisa em espaço de estados (DFID, A*) são \emph{completos} e \emph{óptimos}: se existir uma solução, encontram-na sempre. Os algoritmos de optimização (SA, GA) são \emph{incompletos}: podem não encontrar a solução mesmo que ela exista. Em contrapartida, o SA tem um desempenho muito bom em puzzles fáceis e médios, com tempos inferiores aos do Prolog (overhead de interpretação).

O A* é menos eficiente que o DFID no puzzle Difícil porque a heurística adoptada não distingue estados ao mesmo nível de profundidade, fazendo a fila de prioridade crescer excessivamente. O DFID, com a sua natureza DFS, mantém um número constante de nós em memória, o que o torna mais adequado para puzzles muito restringidos.

\section{Conclusões}
\label{sec:conclusoes}

Este trabalho implementou e comparou quatro abordagens algorítmicas para a resolução do Sudoku clássico e do Killer Sudoku. As principais conclusões são:

\begin{enumerate}
    \item \textbf{DFID com MRV} é o solucionador mais robusto e fiável: resolve todos os puzzles testados, com uso mínimo de memória e tempos razoáveis. A heurística MRV reduz o espaço de busca de forma tão eficaz que o DFID degenera numa DFS com verificação de restrições.

    \item \textbf{A*} partilha a completude do DFID, mas a heurística $h = $ células vazias não distingue estados ao mesmo nível, tornando-o equivalente a uma pesquisa em largura guiada por restrições. Em puzzles com muitos estados alternativos, o A* expande significativamente mais nós que o DFID.

    \item \textbf{SA} é eficaz para puzzles de dificuldade baixa a média, com tempos muito competitivos. A sua principal limitação é o bloqueio em mínimos locais de custo 2, que o operador intra-caixa não consegue resolver directamente a baixa temperatura. A extensão para Killer Sudoku é elegante: basta adicionar o custo de gaiola à função de custo.

    \item \textbf{GA} sofre de convergência prematura e não encontrou soluções completas nos testes realizados. O operador de cruzamento preservador de caixas mantém a validade das restrições de caixa, mas não diversifica suficientemente a população para escapar de mínimos locais profundos.

    \item A extensão para \textbf{Killer Sudoku} foi implementada com sucesso em ambas as linguagens. Em Prolog, a poda por \texttt{partial\_cage\_ok} é muito eficaz, reduzindo o espaço de busca mesmo sem \emph{givens} iniciais. Em Python, a adição do custo de gaiola à função de custo é directa e transparente para os algoritmos SA e GA.
\end{enumerate}

Como trabalho futuro, seria interessante explorar operadores inter-caixa controlados (que preservem as restrições de caixa mas permitam mais mobilidade), técnicas de \emph{reheating} no SA para escapar de mínimos locais profundos, e operadores de cruzamento mais diversificados no GA (por exemplo, cruzamento a nível de linha/coluna em vez de caixa).
